\documentclass[11pt]{article}
\usepackage[utf8]{inputenc}
\usepackage[T1]{fontenc}
\usepackage{amsmath,amssymb}
\usepackage{booktabs}
\usepackage{graphicx}
\usepackage[colorlinks=true,linkcolor=blue,citecolor=blue]{hyperref}
\usepackage[margin=1in]{geometry}
\usepackage{natbib}
\setlength{\parindent}{0pt}
\setlength{\parskip}{6pt}
\providecommand{\tightlist}{}

\title{A Resistive Formalism for Compositional Bottlenecks: \\
A Closed-Form Account of Judgment Correction \textit{(working draft)}}

% author affiliations per D 2026-11-24
\author{
  \textbf{Yifeng Xiao} \thanks{Corresponding author. Email: xiaoyifeng25@mails.ucas.ac.cn}\\
  School of Fundamental Physics and Mathematical Sciences, Hangzhou Institute for Advanced Study, UCAS, Hangzhou 310024, China\\
  University of Chinese Academy of Sciences, Beijing 100049, China\\
  Institute of Theoretical Physics, Chinese Academy of Sciences, Beijing 100190, China
}

\date{\today}
\begin{document}
\maketitle

\section*{Status note}
Working skeleton. Section~2 (the resistive formalism) is the
domain-independent core and is written out. Section~3 instantiates it on GPU
microarchitecture (existing \#53 data). Section~4 is reserved for a physics
instantiation (domain pending decision). Sections~5--6 are stubs.

\section{Introduction}
% TODO: the problem -- where is the bottleneck in a compositional judgment?
% TODO: why a formalism rather than a domain-specific model
% TODO: preview the unifying claim -- one symbol system exports
%       bottleneck + tree + judge simultaneously

\section{The Resistive Formalism}
\label{sec:formalism}

A compositional system is a set of components $i$, each carrying a workload
$f_i$ processed at a rate $r_i$. The time component $i$ occupies is
\begin{equation}
  t_i = \frac{f_i}{r_i}.
  \label{eq:time}
\end{equation}
Two compositions close the system:
\begin{itemize}
  \item \textbf{Series} (sequential phases): $T = \sum_i t_i$;
  \item \textbf{Parallel} (concurrent resources within a phase):
        $t_{\mathrm{phase}} = \max_j\,(f_j/r_j)$.
\end{itemize}
The \emph{bottleneck} is the component that sets the pace:
\begin{equation}
  b^{*} = \operatorname*{argmax}_{j}\; \frac{f_j}{r_j}.
  \label{eq:bottleneck}
\end{equation}
Equations~\eqref{eq:time}--\eqref{eq:bottleneck} are the abstract circuit
law: series adds, parallel is dominated by the largest term, and the
bottleneck is the argmax leaf. One rule ($t=f/r$) plus two aggregations
($\Sigma$, $\max$) covers any two- or three-level structure; the flat model
is the one-level special case, and a branch/segment tree is its two-level
generalization.

\paragraph{Judge as correction.}
A \emph{surface} judgment reads the rates at face value (assumed $r_j$); a
\emph{deep} judgment corrects them (measured $r_j$). The judge score is the
surface-vs-deep correction, decomposed into a \emph{flip} (the bottleneck
identity changes, $b^{*}_{\mathrm{surf}} \neq b^{*}_{\mathrm{deep}}$) and a
\emph{shift} (the ranking within the surviving bottleneck changes). The
ground truth the judge targets is exactly the argmax leaf of
Eq.~\eqref{eq:bottleneck}; the score is the correction distance to it. This
makes the judge a property of the formalism, not an ad-hoc labeler.

\section{Instantiation I: GPU Microarchitecture}
\label{sec:gpu}
The flat model applied to four GPU platforms (workload = a fixed inference
trace; rates = per-component throughput). Table~\ref{tab:gpu} reports the
$2\times2$ tree decomposition (branch-level vs within-branch error).

\begin{table}[h]
\centering
\begin{tabular}{lccc}
\toprule
Platform & rate ratio & within-branch & branch-level \\
\midrule
T5000 & 9.9   & 0.290 & 0.181 \\
A100  & 44.3  & 0.114 & 0.487 \\
5090  & 51.9  & 0.191 & 0.399 \\
4090  & 287.4 & 0.112 & 0.495 \\
\bottomrule
\end{tabular}
\caption{$2\times2$ tree decomposition. Mild-rate platform (T5000):
within-branch dominates; extreme-rate platforms: branch-level dominates.}
\label{tab:gpu}
\end{table}

% TODO: fold in the 1D main result (segment-level flip 0.434/0.556/0.527/0.563)
%       and the flat negative control (both-correct = 1.0)

\section{Instantiation II: Physics (domain TBD)}
\label{sec:physics}
The formalism is domain-agnostic. Candidate physics instantiations:
reaction kinetics (rate-limiting step $=$ argmax of inverse rates), transport
(series resistance / parallel conductance), critical paths.
\textbf{Pending domain decision.}

\section{Discussion: Transferability of the Formalism}
% TODO: what transfers across domains, what does not; predictive content

\section{Conclusion}
% TODO

\section*{AI Use Disclosure}
Per arXiv policy on generative AI tools: no AI system is listed as an author
of this work. Generative AI language tools were used in the drafting and
verification assistance stages; all authors take full responsibility for the
content. All quantitative results are produced by the closed-form computation
scripts referenced in the text and are reproducible from the cited commit
anchors.

\end{document}
